% TOP_PROBLEM: {"problemId": "problem.zilber-s-quasiminimality-conjecture-for-the-complex-exponential-field", "canonicalTitle": "Zilber's Quasiminimality Conjecture for the Complex Exponential Field", "releaseRank": 475, "primaryDomain": "logic_foundations_set_theory", "primaryDomainLabel": "Logic, foundations & set theory", "displayStatus": "open", "releaseStatus": "open"}
% !TeX program = lualatex
% ATTEMPT_STATUS: unresolved
% Model: GPT-6 Astra Ultra; continuation dated 27 September 2026.
\documentclass[11pt]{article}
\usepackage[margin=25mm]{geometry}
\usepackage{fontspec}
\setmainfont{Latin Modern Roman}
\usepackage{amsmath,amssymb,mathtools}
\usepackage{unicode-math}
\setmathfont{Latin Modern Math}
\usepackage{xurl,hyperref}
\hypersetup{colorlinks=true,urlcolor=blue}
\setcounter{secnumdepth}{0}
\setlength{\parindent}{0pt}
\setlength{\parskip}{5pt}
\setlength{\emergencystretch}{3em}
\begin{document}
\section*{\texorpdfstring{Zilber's Quasiminimality Conjecture for the Complex Exponential Field}{Zilber's Quasiminimality Conjecture for the Complex Exponential Field}}\label{475}
\subsection{Definitions and mathematical statement}
The complex exponential field is the first-order structure
\[
 {\mathbb{C}}_{\exp}=({\mathbb{C}};+ ,\cdot,0,1,\exp),\qquad \exp(z)=e^z.
\]
A subset $D\subseteq{\mathbb{C}}$ is definable with parameters if some first-order formula $\varphi(x,a_1,\ldots,a_m)$ in this language, with finitely many parameters from ${\mathbb{C}}$, satisfies $D=\{z:{\mathbb{C}}_{\exp}\models\varphi(z,a_1,\ldots,a_m)\}$. Quantifiers range over all complex numbers.

Quasiminimality asserts that every such one-variable definable set is countable or has countable complement. Countable includes finite. This is not the stronger claim that every definable set is finite or cofinite, nor a claim about arbitrary analytic subsets or arbitrary subsets of ${\mathbb{C}}$. Parameters are allowed, which is part of the requested scope.
\par\smallskip\textit{Formulation and concept references:} \href{https://arxiv.org/pdf/2306.14562v1}{[S1]}.

\subsection{Short English statement}
Is every one-variable set definable using complex arithmetic and exponentiation either countable or the complement of a countable set?
\subsection{Research attempt}

\paragraph{Scope and source audit (27 September 2026).}
The desired dichotomy concerns arbitrary first-order formulas, including
alternations of quantifiers over all of $\mathbb C$, and permits any finite
tuple of complex parameters. It is substantially stronger than the same
statement for zeros of one exponential polynomial. The supplied Wilkie
paper [S1] gives a conditional analytic-continuation route, not a completed
proof for $\mathbb C_{\exp}$. In particular, its Theorem 1 assumes its
Analytic Continuation Property. We do not assume that property below.
The primary-source audit also distinguishes the quasiminimality theorem
for complex powers [E1] from the exponential-field conjecture: the former
uses a weaker structure and does not settle the latter.

This attempt proves the dichotomy for three explicit fragments, derives
a precise common-exceptional-set reformulation of the full problem,
and identifies where unrestricted existential projection defeats the
elementary analytic argument. The arguments are elementary reconstructions,
not claims of new special cases. No identification with a constructed
pseudo-exponential field, and no transcendence conjecture, is assumed.

\paragraph{The countable--cocountable algebra.}
Let $\mathcal A$ be the family of subsets of $\mathbb C$ that are
countable or have countable complement. Countable means at most countable.
This family is closed under complements and countable unions: if one
member of a union is cocountable, so is the union; otherwise the union is
countable. De Morgan's law gives closure under countable intersections.
These assertions use the ordinary set-theoretic countable-union principle.
They do not assert that every countable union is first-order definable.

A useful equivalent description is that any sequence $D_j\in\mathcal A$
has a countable exceptional set $B$ outside which the membership pattern
in all the $D_j$ is constant. Indeed, take the union of $D_j$ when it is
countable and of $\mathbb C\setminus D_j$ otherwise. Each summand is
countable. There cannot be a set whose two sides are both countable,
because $\mathbb C$ is uncountable.

\paragraph{A complete proof for quantifier-free formulas.}
\textbf{Proposition 475.1.} Every subset of $\mathbb C$ defined by a
quantifier-free formula in the stated language, with arbitrary finitely
many complex parameters, belongs to $\mathcal A$.

\emph{Proof.}
Every one-variable term denotes an entire function. This follows by
induction on terms: the variable and constants are entire, sums and
products preserve entire functions, and so does composition with $\exp$.
An atomic equality $t(x)=s(x)$ is the zero set of the entire function
$f=t-s$. If $f$ is identically zero this is all of $\mathbb C$. Otherwise
the identity theorem makes every zero isolated.

For completeness, an isolated subset of $\mathbb C$ is countable.
Choose the countable basis of discs with rational real and imaginary
parts of the centre and positive rational radii. For each isolated point,
choose a basis disc meeting the subset only at that point. Assign the
first such disc in an enumeration. Different points receive different
discs, giving an injection into $\mathbb N$. Thus an atomic equality is
countable or all of $\mathbb C$, and its negation is in $\mathcal A$.
A quantifier-free formula is a finite Boolean combination of atomic
equalities, so the preceding closure properties finish the proof.
\hfill$\square$

The proof allows nested exponentials such as
$\exp(\exp(x+a))$; no restriction on exponential depth was used. It is
not an effective identity test for entire functions represented by terms.
The mathematical alternative ``identically zero or isolated zeros''
does not provide an algorithm deciding the alternative.

The language printed above has no division symbol. If one works on the
natural domain of rational expressions in such terms, denominators can
be cleared while retaining explicit nonvanishing conditions. Each such
condition is still a quantifier-free equality or inequality of entire
terms. Omitting the domain conditions could introduce erroneous points.

\paragraph{Infinite definable sets are already necessary.}
The set
\[
 H=\{z:\exp(z)=1\}=2\pi i\mathbb Z
\]
is countably infinite and has infinite complement. Thus strong minimality
(finite or cofinite definable subsets) fails even in this elementary
fragment. With the allowed parameter $\tau=2\pi i$, the formula
$\exp(\tau x)=1$ defines precisely $\mathbb Z$ as a subset of $\mathbb C$.
This does not contradict quasiminimality: a definable countably infinite
set and its induced arithmetic are permitted by the proposed dichotomy.

Nor may one silently use the real line as a definable set. Both
$\mathbb R$ and $\mathbb C\setminus\mathbb R$ are uncountable, so a
definition of $\mathbb R$ in this structure, even with parameters, would
disprove the target. Taking a real or imaginary part in a mathematical
description is not automatically an operation in its first-order language.
For the same reason the analytic condition $|x|<1$ cannot simply be added
to a formula while claiming to remain in the original structure.

\paragraph{An unrestricted polynomial witness can be eliminated.}
There is a fragment with a quantifier genuinely ranging over all of
$\mathbb C$ that the same method handles completely. Let
\[
 P(x,y)=a_d(x)y^d+\cdots+a_1(x)y+a_0(x),
\]
where each $a_j$ is an entire term in $x$ and the fixed parameters.
\textbf{Proposition 475.2.} The set defined by $\exists y\ P(x,y)=0$
belongs to $\mathcal A$. In fact it is defined by
\[
 \left(\bigvee_{j=1}^{d}a_j(x)\ne0\right)
 \quad\vee\quad
 \left(\bigwedge_{j=0}^{d}a_j(x)=0\right).
\]

\emph{Proof.}
For each fixed $x$, if at least one coefficient with positive index is
nonzero, $P(x,\cdot)$ is a nonconstant complex polynomial and has a
root. If all those coefficients vanish, it has a root exactly when the
constant coefficient vanishes, in which case every $y$ is a root.
These are all possibilities. Proposition 475.1 applies to the displayed
formula.\hfill$\square$

The proof does not apply to arbitrary systems by eliminating one equation
at a time: a witness for one equation need not also solve another. Nor
does it cover $\exp(y)$ occurring in the witness variable. As a simple
transcendental instance that is nevertheless explicit,
\[
 \{x:\exists y\ \exp(y)=x\}=\mathbb C\setminus\{0\}.
\]
The exclusion of zero is immediate. For $x\ne0$, choose an argument of
$x$ and take $y=\log|x|+i\arg x$. This is an existence proof, not a
claim that a global definable choice of logarithm has been constructed.

\paragraph{Countable witness sets and a quantified fragment.}
\textbf{Proposition 475.3.} Let $B_1,\ldots,B_m$ be fixed countable
subsets of $\mathbb C$. Start with any quantifier-free formula
$\theta(x,y_1,\ldots,y_m,a)$ and quantify its $y_j$ variables in any
order using $\exists y_j\in B_j$ or $\forall y_j\in B_j$.
The resulting set of $x$ belongs to $\mathcal A$.

\emph{Proof.}
For each fixed tuple of witnesses, Proposition 475.1 applies, since the
witness entries become finitely many additional complex parameters.
An existential quantifier over a countable set takes a countable union;
a universal quantifier takes a countable intersection. Induction outward
through the finite quantifier string proves the assertion. Empty witness
sets cause no problem: the respective results are empty or all of
$\mathbb C$.\hfill$\square$

If the $B_j$ are definable, these restricted quantifiers can of course
be expressed by ordinary formulas. For example $H=\ker(\exp)$ is a
parameter-free definable countable witness set. The proposition covers
arbitrary finite alternations over that fixed kernel and allows arbitrary
nested exponential terms in the matrix. It also covers fixed countable
fibres $\{y:\exp(y)=c\}$, when parameters $c$ are fixed.

An important qualification is the word \emph{fixed}. For the relation
$\exp(y)=x$, each fibre over a nonzero $x$ is countable, but the union
of all fibres is $\mathbb C$. The preceding proof does not turn
an $x$-dependent countable family of witnesses into a fixed countable
index set. A local choice of logarithm changes by $2\pi i\mathbb Z$,
and the choice of that logarithm is precisely extra information.

\paragraph{One common exceptional set: an exact reformulation.}
Let $A\subset\mathbb C$ be countable, and let $\operatorname{Form}_1(A)$
denote all formulas with one free variable and finitely many parameters
from $A$. There are only countably many such formulas: the language is
countable, every formula is finite, and finite tuples from $A$ form a
countable set.

\textbf{Proposition 475.4.} The conjecture is equivalent to the following
assertion for every countable $A$: there exists a countable set $C_A$
such that any two points outside $C_A$ satisfy exactly the same formulas
in $\operatorname{Form}_1(A)$.

\emph{Proof.}
Assume quasiminimality. For each formula, take its countable side, either
its realization set or the complement. Their countable union is $C_A$.
Outside it, each formula has a constant truth value, so any two such
points have the same full one-variable type over $A$.

Conversely, take a formula with a finite parameter set $A$ and a set
$C_A$ as described. Its truth value is constant on the nonempty set
$\mathbb C\setminus C_A$. The defined set is therefore contained in
$C_A$, or its complement is contained in $C_A$. This is exactly the
desired dichotomy.\hfill$\square$

This is a reformulation, not a construction of $C_A$. In particular it
does not say that two points of the same type are interchanged by an
automorphism of the given, nonsaturated structure. Such an automorphism
statement would require its own homogeneity argument.

One can explicitly construct a countable field $F_A$ containing
$A\cup\{2\pi i\}$ and closed under algebraic closure inside $\mathbb C$,
exponentiation, and all logarithms of its nonzero elements. Start with
the countable field generated by $A\cup\{2\pi i\}$, and repeatedly
adjoin its algebraic elements, its exponential images, and the fibres
of $\exp$ over its nonzero elements; then take the field generated by
these elements. Each step remains countable: a countable field has
countably many polynomials with finitely many roots each, and each
nonempty exponential fibre is a coset of $2\pi i\mathbb Z$.
The union of the stages has the claimed closure properties.

What is missing is a proof that all points outside this, or some larger
countable set, have the same type. Closure under field operations,
exponentiation, and logarithms does not control arbitrary quantified
systems. Choosing $C_A=F_A$ without proving this extra assertion would
simply conceal the main difficulty.

\paragraph{Where the analytic projection attack stops.}
Suppose a system of entire equations $F(x,y)=0$, with
$y\in\mathbb C^r$ and $F$ taking values in $\mathbb C^r$, has a solution
$(x_0,y_0)$ with $\det(\partial F/\partial y)(x_0,y_0)\ne0$.
The holomorphic implicit function theorem supplies an open disc $U$
around $x_0$ and a holomorphic witness $y(x)$ on $U$.
Consequently its existential projection contains $U$.
Any finitely many strict nonvanishing conditions true at $(x_0,y_0)$
remain true after shrinking $U$, by continuity.

This is genuine local information, but it is not cocountability. Even
an open subset of $\mathbb C$ can have uncountable complement. The unit
disc is a topological illustration of that logical gap, not an asserted
exponential-field-definable counterexample. Singular witnesses present
an additional problem, since the displayed Jacobian condition can fail.

Attempting to move the local witness along paths requires continuation
and control of the exceptional set. A local logarithm demonstrates the
distinction between local solvability and global single-valued choice.
If a holomorphic $g$ on $\mathbb C^*$ satisfied $e^{g(z)}=z$, then
differentiation would give $g'(z)=1/z$. Integrating around the unit
circle yields $0=2\pi i$, impossible. Local logarithms can be continued
along paths, but the resulting value depends on winding. General
implicit exponential systems require substantially more than this
well-understood example, and one cannot assume their witnesses extend
along every desired path.

Wilkie's Analytic Continuation Property supplies a carefully qualified
continuation statement along generic paths; his Theorem 1 then obtains
quasiminimality using a back-and-forth argument. The elementary lemmas
above neither establish that property nor replace the back-and-forth
step. The source is used to locate the missing global mechanism,
not to import its unproved hypothesis as a theorem about $\mathbb C_{\exp}$.

\paragraph{Adversarial checks and outcome.}
The parameter issue was tested at each step: specializing witnesses
creates only finitely many parameters per slice, whereas the common
exceptional-set lemma explicitly takes the union over countably many
formulas. The countable language alone cannot force quasiminimality:
adding one unary predicate for $\mathbb R$ still gives a countable
language and immediately gives a set with two uncountable sides.
Similarly, countability of a generated subfield cannot imply uniformity
of all types outside it.

What has been proved is the full dichotomy for quantifier-free formulas,
for a single polynomial witness with entire-term coefficients, and for
arbitrary finite strings of quantifiers over fixed countable witness
sets. Proposition 475.4 gives an exact target for extending these results.
The first unresolved step is uniform control of unrestricted complex
witnesses and their continuation, strong enough to handle every formula
over a fixed countable parameter set. No proof or counterexample for
the full conjecture is supplied.

\subsection{Sources}
\begin{itemize}
\item[S1]A. J. Wilkie, "Analytic continuation and Zilber's quasiminimality conjecture," arXiv:2306.14562 (26 June 2023).
\url{https://arxiv.org/pdf/2306.14562v1}.
\textit{Formulation source.}
\item[E1]Francesco Gallinaro and Jonathan Kirby, \emph{Quasiminimality
of complex powers}, Forum of Mathematics, Sigma 12 (2024), e123.
\url{https://arxiv.org/abs/2304.06450};
\url{https://doi.org/10.1017/fms.2024.82}.
\textit{Primary abstract and accepted-manuscript introduction inspected
27 September 2026. The powers structure is distinguished from the full
complex exponential field.}
\item[E2]Alex J. Wilkie, full-text HTML of the supplied source,
\url{https://arxiv.org/html/2306.14562v1}.
\textit{Definitions 1--3 and conditional Theorem 1 inspected
27 September 2026; the Analytic Continuation Property is not treated
as proved for the structure needed here.}
\end{itemize}
\end{document}
